\documentclass[11pt]{article}
\usepackage[a4paper,margin=25mm]{geometry}
\usepackage[T1]{fontenc}
\usepackage[utf8]{inputenc}
\usepackage{textcomp}
\usepackage{amsmath,amssymb}
\usepackage{xcolor}
\usepackage{booktabs,longtable,array}
\usepackage{enumitem}
\usepackage{listings}
\usepackage[hidelinks]{hyperref}

\setlength{\parindent}{0pt}
\setlength{\parskip}{0.55em}
\setlist{nosep,leftmargin=*}
\emergencystretch=3em

\lstdefinestyle{logic}{
  basicstyle=\ttfamily\small,
  columns=fullflexible,
  keepspaces=true,
  breaklines=true,
  frame=single,
  framerule=0.3pt,
  rulecolor=\color{black!35},
  xleftmargin=0.5em,
  xrightmargin=0.5em,
  aboveskip=0.7em,
  belowskip=0.7em,
  showstringspaces=false
}

\title{M12x --- Cell-by-cell logical inspection (18 cells)}
\author{}
\date{}
\begin{document}
\maketitle

\textbf{Experiment ID:} M12x (Milestone 1 Part 2 expanded)\\
\textbf{Status:} Stage-3 qualitative inspection \textbf{in progress} (12/18 written)\\
\textbf{Date started:} 2026-07-20\\
\textbf{Parent record:} \href{M1.2-expanded.md}{\texttt{M1.2-expanded.md}}\\
\textbf{Artefacts:} \texttt{causal/outputs/aba\_learning/grid/M12x\_\{ecai2024,aamas2025\}/cells/\{dgp\}\_\_target-\{t\}/}\\
\textbf{Side-car matrix:} \texttt{causal/outputs/aba\_learning/grid/M12x\_summary.md}\\
\textbf{Per-cell detector reports:} \texttt{causal/outputs/aba\_learning/grid/M12x\_cell\_reports/}\\
\textbf{Fixtures:} \texttt{causal/experiments/handcrafted\_m12x.py}\\
\textbf{Catalogue:} \texttt{docs/research/milestone\_plans/milestone1\_part2/m12x\_units\_reference.tex}\\
\textbf{Grid run:} 2026-07-20 (18/18 cells \texttt{solved}; summary regenerated same day)

This document inspects each of the 18 \texttt{(arm \(\times\) fixture \(\times\) target)} cells. For every completed cell it records: the locked setup and \(\mathcal{H}_t^\star\); a clear reading of how ABA Learning produced the learned framework; and whether and why that framework matches or fails the intended hypothesis and semantic success description.

This is \textbf{not} Russo-style Causal ABA. The object of study is unguided ABA Learning recovery of mechanism-aligned target rules from categorical tabular background knowledge and positive/negative examples.

\textbf{Reading conventions.} Learned \(\delta\) is the new intensional material in \texttt{bk.sol.aba} after variable-canonical background knowledge subtraction. Exact match means the learned target rules equal \(\mathcal{H}_t^\star\) after normalisation. ASP \texttt{pos\_covered} / \texttt{neg\_rejected} fractions are coverage facts from the side-car; they are not a substitute for intensional match. Acceptance remains inspection-first against the Approach semantic description.

\medskip\hrule\medskip

\section{Inspection checklist (18)}

{\small
\setlength{\tabcolsep}{3pt}
\begin{longtable}{@{}>{\centering\arraybackslash}p{0.035\textwidth}p{0.075\textwidth}p{0.235\textwidth}>{\centering\arraybackslash}p{0.055\textwidth}>{\centering\arraybackslash}p{0.105\textwidth}p{0.355\textwidth}@{}}
\toprule
\# & Arm & Fixture & Target & Exact \(\mathcal{H}^\star\)? & Short verdict \\
\midrule
\endfirsthead
\toprule
\# & Arm & Fixture & Target & Exact \(\mathcal{H}^\star\)? & Short verdict \\
\midrule
\endhead
1 & ECAI & \texttt{m12\_u1\_separator\_copy} & \(x_2\) & N & Assumption-mediated via isolated \(x_0\); parent only in contraries \\
2 & ECAI & \texttt{m12\_u2\_collider\_min} & \(x_2\) & N & Elegant assumption-mediated both-parents-nonzero (not the locked 4-clause form) \\
3 & ECAI & \texttt{m12\_u3\_collider\_max} & \(x_2\) & N & \(x_0\)-case OR: pure unary on \(x_0\in\{1,2\}\); assumption only on \(x_0=0\) slice \\
4 & ECAI & \texttt{m12\_u4\_fork\_asymmetric} & \(x_1\) & Y & Exact \(\mathcal{H}^\star\): two unary parent rules; sibling unused \\
5 & ECAI & \texttt{m12\_u4\_fork\_asymmetric} & \(x_2\) & Y & Exact \(\mathcal{H}^\star\): two unary parent rules; sibling unused (mirror of cell 4) \\
6 & ECAI & \texttt{m12\_u5\_chain\_curated} & \(x_2\) & N & Ancestor \(x_0\) case + assumptions; parent only in contraries \\
7 & ECAI & \texttt{m12\_u6\_g1\_min\_diff} & \(x_2\) & --- & \emph{pending} \\
8 & ECAI & \texttt{m12\_u6\_g1\_min\_diff} & \(x_3\) & --- & \emph{pending} \\
9 & ECAI & \texttt{m12\_u7\_diamond\_noisy} & \(x_3\) & --- & \emph{pending} \\
10 & AAMAS & \texttt{m12\_u1\_separator\_copy} & \(x_2\) & N & Six specialised conjunctions; isolated \(x_0\) retained beside parent \\
11 & AAMAS & \texttt{m12\_u2\_collider\_min} & \(x_2\) & Y & Exact \(\mathcal{H}^\star\): four both-parent conjunctions, no assumptions \\
12 & AAMAS & \texttt{m12\_u3\_collider\_max} & \(x_2\) & N & Eight pairwise positive-cell conjunctions; not the locked unary OR \\
13 & AAMAS & \texttt{m12\_u4\_fork\_asymmetric} & \(x_1\) & N & Parent + sibling conjunctions; cites distractor \(x_2\) \\
14 & AAMAS & \texttt{m12\_u4\_fork\_asymmetric} & \(x_2\) & N & Parent + sibling conjunctions; cites distractor \(x_1\) \\
15 & AAMAS & \texttt{m12\_u5\_chain\_curated} & \(x_2\) & N & Parent + grandparent conjunctions; ancestor retained \\
16 & AAMAS & \texttt{m12\_u6\_g1\_min\_diff} & \(x_2\) & --- & \emph{pending} \\
17 & AAMAS & \texttt{m12\_u6\_g1\_min\_diff} & \(x_3\) & --- & \emph{pending} \\
18 & AAMAS & \texttt{m12\_u7\_diamond\_noisy} & \(x_3\) & --- & \emph{pending} \\
\bottomrule
\end{longtable}
}

\medskip\hrule\medskip

\section{Shared regime (locked)}

\begin{itemize}
  \item Alphabet \(K=\{0,1,2\}\); nonzero-positive labels \(E^+=\{t\neq 0\}\), \(E^-=\{t=0\}\) unless a unit card overrides.
  \item Background knowledge: exact-value predicates \texttt{xi\_val\_v} for every non-target column (descendants kept as distractors).
  \item Configs: ECAI (\texttt{configs/ecai2024\_config.pl}); AAMAS (\texttt{configs/aamas2025\_config.pl}).
  \item Reference hypotheses \(\mathcal{H}_t^\star\) and structural sets are as in the unit catalogue and \texttt{m12x\_summary.py} look-up tables.
\end{itemize}

\textbf{Effective options (confirmed in cell traces):}

\begin{itemize}
  \item \textbf{ECAI:} brave learning; \texttt{folding\_mode(nd)}; \texttt{folding\_selection(any)}; \texttt{folding\_space(all)}; \texttt{asm\_intro(relto)}; \texttt{post\_folding\_test\_entailment(true)}.
  \item \textbf{AAMAS:} brave learning; greedy folding; published AAMAS options consulted verbatim from \texttt{configs/aamas2025\_config.pl}.
\end{itemize}

\medskip\hrule\medskip

\section{Batch A --- ECAI (\texttt{M12x\_ecai2024})}

\subsection{1. ECAI \(\times\) \texttt{m12\_u1\_separator\_copy} \(\times\) \(t=x_2\)}

\textbf{Cell directory:} \texttt{M12x\_ecai2024/cells/m12\_u1\_separator\_copy\_\_target-x2/}\\
\textbf{Runner outcome:} \texttt{solved}\\
\textbf{ASP coverage:} positives \texttt{6/6}, negatives \texttt{3/3}

\subsubsection{Setup}

The graph is a separator \(x_1\to x_2\) with an isolated variable \(x_0\). The mechanism sets \(x_2\) equal to \(x_1\) on the full nine-row factorial over \((x_0,x_1)\). Under nonzero-positive labelling the positive examples are rows \(\{2,3,5,6,8,9\}\) and the negative examples are rows \(\{1,4,7\}\). Background knowledge contains exact-value predicates for \(x_0\) and \(x_1\); the target column \(x_2\) is excluded. The ancestors of \(x_2\) are \(\{x_1\}\); the isolated variable \(x_0\) is a distractor; there are no descendants. The design role of this unit is to test whether the learner recovers a unary parent rule and avoids citing the isolated distractor.

\subsubsection{Intended hypothesis \(\mathcal{H}_{x_2}^\star\)}

\begin{lstlisting}[style=logic,language=Prolog]
x2(A) :- x1_val_1(A).
x2(A) :- x1_val_2(A).
\end{lstlisting}

These two clauses characterise \(x_2\neq 0\) by the parent \(x_1\) being nonzero. They must not cite the isolated distractor \(x_0\).

\subsubsection{What happened}

This cell uses the ECAI published configuration on unit U1 with learning target \(x_2\). The learner returned a solved framework with full ASP coverage of the examples. That framework does not match \(\mathcal{H}_{x_2}^\star\) intensionally. The target rules cite the isolated distractor \(x_0\) together with assumptions, and the true parent \(x_1\) appears only in the contrary conditions.

\subsubsection{What the learner did}

Learning begins from rote rules that cover each positive example by its sample identifier:

\begin{lstlisting}[style=logic,language=Prolog]
x2(A) :- A=2.
x2(A) :- A=3.
x2(A) :- A=5.
x2(A) :- A=6.
x2(A) :- A=8.
x2(A) :- A=9.
\end{lstlisting}

Under non-deterministic folding with selection \texttt{any} and folding space \texttt{all}, the first successful fold for a representative positive in each \(x_0\)-slice replaces that identifier by a value predicate of the isolated column \(x_0\). The first fold yields a body containing only \texttt{x0\_val\_0}; later folds yield only \texttt{x0\_val\_1} and only \texttt{x0\_val\_2}:

\begin{lstlisting}[style=logic,language=Prolog]
x2(A) :- x0_val_0(A).   % from A=2; overgeneralises within the x0=0 slice
x2(A) :- x0_val_1(A).   % from A=5
x2(A) :- x0_val_2(A).   % from A=8
\end{lstlisting}

Each of those bodies overgeneralises within its slice, because the negative example that shares the same \(x_0\) value would also be covered. Entailment of the labelled example set therefore fails, and assumption introduction relative to the failing body adds a fresh assumption \texttt{alpha\_1}, \texttt{alpha\_2}, or \texttt{alpha\_3}. The remaining positive examples in the same slice are then deleted as subsumed. The repaired target rules are:

\begin{lstlisting}[style=logic,language=Prolog]
x2(A) :- alpha_1(A), x0_val_0(A).
x2(A) :- alpha_2(A), x0_val_1(A).
x2(A) :- alpha_3(A), x0_val_2(A).
\end{lstlisting}

Contrary learning is driven by the three negative examples, which are exactly the rows with \(x_1=0\). For each contrary, folding first tries the matching \(x_0\) literal already used in the corresponding target rule. That attempt is rejected, because the engine cannot introduce an assumption already tied to that body. Folding then succeeds with \texttt{x1\_val\_0}:

\begin{lstlisting}[style=logic,language=Prolog]
c_alpha_1(A) :- x1_val_0(A).
c_alpha_2(A) :- x1_val_0(A).
c_alpha_3(A) :- x1_val_0(A).
\end{lstlisting}

The accepted solution therefore has three assumption-mediated target rules keyed by \(x_0\), and three identical contrary rules that block those assumptions whenever \(x_1=0\). Body-scope variables of the target rules are \(\{x_0\}\). The parent \(x_1\) appears only in the framework-scope contrary conditions.

\subsubsection{Match assessment}

The learned target rules are not an exact match to \(\mathcal{H}_{x_2}^\star\).

On the complete partition of the table by \(x_0\), an assumption remains undefeated on a row in that slice if and only if \texttt{x1\_val\_0} fails, which is exactly when \(x_1\in\{1,2\}\). Taken together, the three rules therefore cover the six positives and reject the three negatives. In that limited extensional sense the solution agrees with the nonzero-positive copy semantics of U1.

Intensionally the solution is misaligned with the intended hypothesis. The learner folded the isolated distractor \(x_0\) before the true parent \(x_1\). It repaired the resulting overgeneralisation by introducing assumptions. It recovered the true separator only as contrary conditions on \(x_1\). The intended hypothesis cites \(x_1\) directly in ordinary target rules and does not mention \(x_0\) at all.

\textbf{Short verdict.} The cell is solved and extensionally adequate, but intensionally it encodes the parent concept through assumptions keyed by the isolated distractor, with the true parent appearing only in the contraries.

\medskip\hrule\medskip

\subsection{2. ECAI \(\times\) \texttt{m12\_u2\_collider\_min} \(\times\) \(t=x_2\)}

\textbf{Cell directory:} \texttt{M12x\_ecai2024/cells/m12\_u2\_collider\_min\_\_target-x2/}\\
\textbf{Runner outcome:} \texttt{solved}\\
\textbf{ASP coverage:} positives \texttt{4/4}, negatives \texttt{5/5}

\subsubsection{Setup}

Unit U2 is a collider \(x_0\to x_2\leftarrow x_1\) with mechanism \(x_2:=\min(x_0,x_1)\) on the full nine-row factorial over \((x_0,x_1)\). Under nonzero-positive labelling, \(x_2\neq 0\) if and only if both parents are nonzero. The positive examples are therefore rows \(\{5,6,8,9\}\) and the negative examples are rows \(\{1,2,3,4,7\}\). Background knowledge contains exact-value predicates for both parents; the target \(x_2\) is excluded. There are no sibling or descendant distractors beyond the two parents themselves. The design role of this unit is to test recovery of a conjunctive both-parent rule. The locked reference \(\mathcal{H}_{x_2}^\star\) enumerates that concept as four ordinary \texttt{val} conjunctions, one for each nonzero pair of parent values.

\subsubsection{Intended hypothesis \(\mathcal{H}_{x_2}^\star\)}

\begin{lstlisting}[style=logic,language=Prolog]
x2(A) :- x0_val_1(A), x1_val_1(A).
x2(A) :- x0_val_1(A), x1_val_2(A).
x2(A) :- x0_val_2(A), x1_val_1(A).
x2(A) :- x0_val_2(A), x1_val_2(A).
\end{lstlisting}

\subsubsection{What happened}

This cell uses the ECAI published configuration on unit U2 with learning target \(x_2\). The learner returned a solved framework with full ASP coverage. It did not emit the four locked case clauses. Instead it recovered a compact assumption-mediated representation of the same semantic idea: \(x_0\) is nonzero in the target bodies, and \(x_1\) is nonzero through undefeated assumptions whose contraries fire exactly when \(x_1=0\).

\subsubsection{What the learner did}

Learning begins from rote rules for the four positives:

\begin{lstlisting}[style=logic,language=Prolog]
x2(A) :- A=5.
x2(A) :- A=6.
x2(A) :- A=8.
x2(A) :- A=9.
\end{lstlisting}

Under non-deterministic folding with selection \texttt{any} and folding space \texttt{all}, the first fold for \(A=5\) (where \(x_0=1\)) replaces the sample identifier by \texttt{x0\_val\_1} alone:

\begin{lstlisting}[style=logic,language=Prolog]
x2(A) :- x0_val_1(A).
\end{lstlisting}

That body overgeneralises: it would also cover the negative row \(A=4\), where \(x_0=1\) but \(x_1=0\). Entailment of the labelled example set therefore fails, and assumption introduction adds \texttt{alpha\_1}:

\begin{lstlisting}[style=logic,language=Prolog]
x2(A) :- alpha_1(A), x0_val_1(A).
\end{lstlisting}

The remaining positive in the same \(x_0=1\) slice (\(A=6\)) is then deleted as subsumed. The same pattern occurs for the \(x_0=2\) slice: fold \(A=8\) to \texttt{x0\_val\_2}, introduce \texttt{alpha\_2}, and delete \(A=9\) as subsumed:

\begin{lstlisting}[style=logic,language=Prolog]
x2(A) :- alpha_2(A), x0_val_2(A).
\end{lstlisting}

Contrary learning is driven by the negatives that sit in those same \(x_0\) slices, namely \(A=4\) and \(A=7\) (both have \(x_1=0\)). Folding each contrary first tries the matching \(x_0\) literal already used in the corresponding target rule; that attempt is rejected. Folding then succeeds with \texttt{x1\_val\_0}:

\begin{lstlisting}[style=logic,language=Prolog]
c_alpha_1(A) :- x1_val_0(A).
c_alpha_2(A) :- x1_val_0(A).
\end{lstlisting}

In the ASP solution these assumptions are realised by default-negation rules scoped to the same \(x_0\)-slices:

\begin{lstlisting}[style=logic,language=Prolog]
alpha_1(A) :- not c_alpha_1(A), x0_val_1(A).
alpha_2(A) :- not c_alpha_2(A), x0_val_2(A).
\end{lstlisting}

Those clauses are what allow each assumption to hold on its slice precisely when the corresponding contrary fails, that is when \(x_1\neq 0\). Without them, \texttt{alpha\_1} and \texttt{alpha\_2} would never be derived, and the target rules that depend on them would not cover the positives.

Negatives with \(x_0=0\) (rows \(\{1,2,3\}\)) never enter these target bodies, so they need no contrary repair. Body-scope variables of the target rules are \(\{x_0\}\). The second parent \(x_1\) appears in the contrary conditions, and thereby controls whether the assumptions succeed.

\subsubsection{Match assessment}

Relative to the locked case enumeration \(\mathcal{H}_{x_2}^\star\), this is not an exact syntactic match. The learner does not produce the four ordinary both-parent conjunctions.

Judged against the semantic success description for U2, however, the learned framework is a strong and arguably more elegant intensional solution. It characterises nonzero \(x_2\) under \(\min\) by combining an explicit nonzero condition on one parent in the target bodies with an assumption-mediated nonzero condition on the other parent. On an \(x_0\in\{1,2\}\) row, \(x_2\) holds if and only if the assumption is undefeated if and only if \(x_1\neq 0\). That is exactly ``both parents are nonzero,'' expressed with two target rules and the ABA assumption machinery rather than with four case clauses. ASP coverage is complete for that reason.

So the mismatch with \(\mathcal{H}_{x_2}^\star\) should not be read as a failure to recover the intended concept. It is a mismatch with one particular syntactic presentation of that concept. The learned form uses assumptions to represent nonzeroness compactly.

\textbf{Short verdict.} Not an exact match to the locked four-clause \(\mathcal{H}_{x_2}^\star\), but an elegant assumption-mediated encoding of the same both-parents-nonzero concept, with full coverage.

\medskip\hrule\medskip

\subsection{3. ECAI \(\times\) \texttt{m12\_u3\_collider\_max} \(\times\) \(t=x_2\)}

\textbf{Cell directory:} \texttt{M12x\_ecai2024/cells/m12\_u3\_collider\_max\_\_target-x2/}\\
\textbf{Runner outcome:} \texttt{solved}\\
\textbf{ASP coverage:} positives \texttt{8/8}, negatives \texttt{1/1}

\subsubsection{Setup}

Unit U3 is the disjunctive counterpart of U2: the same collider \(x_0\to x_2\leftarrow x_1\), now with mechanism \(x_2:=\max(x_0,x_1)\) on the nine-row factorial. Under nonzero-positive labelling, \(x_2\neq 0\) if and only if at least one parent is nonzero. The positive examples are therefore rows \(\{2,3,4,5,6,7,8,9\}\) and the sole negative is row \(\{1\}\). Background knowledge contains exact-value predicates for both parents; the target \(x_2\) is excluded. The locked reference \(\mathcal{H}_{x_2}^\star\) is the four-clause \texttt{val} expansion of ``at least one parent is nonzero.''

\subsubsection{Intended hypothesis \(\mathcal{H}_{x_2}^\star\)}

\begin{lstlisting}[style=logic,language=Prolog]
x2(A) :- x0_val_1(A).
x2(A) :- x0_val_2(A).
x2(A) :- x1_val_1(A).
x2(A) :- x1_val_2(A).
\end{lstlisting}

\subsubsection{What happened}

This cell uses the ECAI published configuration on unit U3 with learning target \(x_2\). After the initial optimum there are eight rote rules, one for each positive. Non-deterministic folding proceeds by \(x_0\)-slices and yields a mixed solution: pure unary rules for \(x_0\in\{1,2\}\), and an assumption-mediated rule only on the incomplete \(x_0=0\) slice. The learner returns a solved framework with full ASP coverage. The result is not an exact syntactic match to \(\mathcal{H}_{x_2}^\star\), but it is a compact case analysis that still expresses the intended disjunctive concept.

\subsubsection{What the learner did}

Learning begins from rote coverage of the eight positives. The first fold chosen is \(A=2\) (where \(x_0=0\), \(x_1=1\)), which replaces the identifier by \texttt{x0\_val\_0} alone:

\begin{lstlisting}[style=logic,language=Prolog]
x2(A) :- x0_val_0(A).
\end{lstlisting}

That body would also cover the unique negative \(A=1\) (where \(x_0=0\) and \(x_1=0\)), so entailment fails and assumption introduction adds \texttt{alpha\_1}:

\begin{lstlisting}[style=logic,language=Prolog]
x2(A) :- alpha_1(A), x0_val_0(A).
\end{lstlisting}

The other positive in the same slice (\(A=3\)) is then deleted as subsumed. Next, \(A=4\) folds to \texttt{x0\_val\_1}. Under \(\max\), every row with \(x_0=1\) is already positive, so this one-literal rule entails \(\langle E^+,E^-\rangle\) immediately and needs no assumption:

\begin{lstlisting}[style=logic,language=Prolog]
x2(A) :- x0_val_1(A).
\end{lstlisting}

Positives \(A=5\) and \(A=6\) are subsumed. Likewise \(A=7\) folds to a safe unary rule on the \(x_0=2\) slice:

\begin{lstlisting}[style=logic,language=Prolog]
x2(A) :- x0_val_2(A).
\end{lstlisting}

Positives \(A=8\) and \(A=9\) are subsumed. Contrary learning for the negative \(A=1\) first tries \texttt{x0\_val\_0} (rejected), then succeeds with \texttt{x1\_val\_0}:

\begin{lstlisting}[style=logic,language=Prolog]
c_alpha_1(A) :- x1_val_0(A).
\end{lstlisting}

In the ASP solution the assumption is realised by

\begin{lstlisting}[style=logic,language=Prolog]
alpha_1(A) :- not c_alpha_1(A), x0_val_0(A).
\end{lstlisting}

so on the \(x_0=0\) slice the assumption holds precisely when \(x_1\neq 0\). Body-scope variables of the target rules are \(\{x_0\}\); \(x_1\) appears only through the contrary / assumption pair.

\subsubsection{Match assessment}

This is not an exact match to \(\mathcal{H}_{x_2}^\star\). The learner does not emit the four unary parent rules of the locked OR expansion.

Extensionally, the learned theory is

\[ (x_0=0 \land x_1\neq 0)\;\lor\;(x_0=1)\;\lor\;(x_0=2), \]

which is exactly \(x_0\neq 0\) or \(x_1\neq 0\), the nonzero-positive semantics of \(\max\). ASP coverage is complete for that reason.

Intensionally this is an \(x_0\)-first case analysis: the complete nonzero \(x_0\) slices become ordinary unary rules, and only the incomplete \(x_0=0\) slice needs an assumption to encode ``\(x_1\) nonzero.'' Relative to \(\mathcal{H}_{x_2}^\star\), the presentation differs (no free-standing \texttt{x1\_val\_1} / \texttt{x1\_val\_2} target rules), but the recovered concept is the intended disjunction. The pattern is related to ECAI on U2: assumption machinery is used where a one-literal fold would overgeneralise, while folds that are already safe are kept without assumptions.

\textbf{Short verdict.} Not an exact match to the locked four-rule OR, but a solved, fully covering \(x_0\)-case analysis that encodes the same ``at least one parent nonzero'' concept, with assumptions only on the \(x_0=0\) slice.

\medskip\hrule\medskip

\subsection{4. ECAI \(\times\) \texttt{m12\_u4\_fork\_asymmetric} \(\times\) \(t=x_1\)}

\textbf{Cell directory:} \texttt{M12x\_ecai2024/cells/m12\_u4\_fork\_asymmetric\_\_target-x1/}\\
\textbf{Runner outcome:} \texttt{solved}\\
\textbf{ASP coverage:} positives \texttt{2/2}, negatives \texttt{1/1}

\subsubsection{Setup}

Unit U4 is a fork \(x_0\to x_1\), \(x_0\to x_2\) on the three-row source factorial over \(x_0\in\{0,1,2\}\). The asymmetric maps are \(x_1:=2\) if \(x_0\neq 2\) else \(0\), and \(x_2:=2\) if \(x_0\neq 0\) else \(0\). For target \(x_1\), nonzero-positive labelling gives \(E^+=\{1,2\}\) and \(E^-=\{3\}\); on the table, \(x_1\neq 0\) if and only if \(x_0\in\{0,1\}\). Background knowledge includes exact-value predicates for the parent \(x_0\) and the sibling distractor \(x_2\); the target \(x_1\) is excluded. The sibling is correlated with \(x_1\) via \(x_0\) but imperfect as a separator (\texttt{x2\_val\_2} holds on both positive row 2 and negative row 3). The locked reference \(\mathcal{H}_{x_1}^\star\) is the two-clause parent expansion of \(x_0\neq 2\).

\subsubsection{Intended hypothesis \(\mathcal{H}_{x_1}^\star\)}

\begin{lstlisting}[style=logic,language=Prolog]
x1(A) :- x0_val_0(A).
x1(A) :- x0_val_1(A).
\end{lstlisting}

\subsubsection{What happened}

This cell uses the ECAI published configuration on unit U4 with learning target \(x_1\). After the initial optimum there are two rote rules, one for each positive. Non-deterministic folding replaces each identifier by the corresponding parent value in a single step, and both folds entail \(\langle E^+,E^-\rangle\) immediately. The learner returns a solved framework with full ASP coverage, no assumptions, and no citation of the sibling.

\subsubsection{What the learner did}

Learning begins from rote coverage of the two positives:

\begin{lstlisting}[style=logic,language=Prolog]
x1(A) :- A=1.
x1(A) :- A=2.
\end{lstlisting}

The first fold chosen is \(A=1\) (where \(x_0=0\)), which replaces the identifier by \texttt{x0\_val\_0}:

\begin{lstlisting}[style=logic,language=Prolog]
x1(A) :- x0_val_0(A).
\end{lstlisting}

That body covers only the positive row with \(x_0=0\) and does not cover the negative (\(x_0=2\)), so entailment succeeds without assumption introduction. The second fold is \(A=2\) (where \(x_0=1\)), which likewise yields a safe unary parent rule:

\begin{lstlisting}[style=logic,language=Prolog]
x1(A) :- x0_val_1(A).
\end{lstlisting}

Learning then stops with \texttt{nothing to fold}. Although sibling predicates (\texttt{x2\_val\_0}, \texttt{x2\_val\_2}) are present in the background knowledge, neither appears in a target body. Body-scope variables of the target rules are \(\{x_0\}\).

\subsubsection{Match assessment}

This is an exact match to \(\mathcal{H}_{x_1}^\star\).

The two unary parent rules are precisely the locked \texttt{val} expansion of ``\(x_0\neq 2\)'' under nonzero-positive labelling of \(x_1\). ASP coverage is complete, the solution uses no assumptions, and the sibling distractor is not cited. Relative to the unit's semantic success description, the learned framework characterises nonzero \(x_1\) via the true parent without rote sample-id casework.

\textbf{Short verdict.} Exact match to \(\mathcal{H}_{x_1}^\star\): two unary parent rules, no assumptions, sibling unused, full coverage.

\medskip\hrule\medskip

\subsection{5. ECAI \(\times\) \texttt{m12\_u4\_fork\_asymmetric} \(\times\) \(t=x_2\)}

\textbf{Cell directory:} \texttt{M12x\_ecai2024/cells/m12\_u4\_fork\_asymmetric\_\_target-x2/}\\
\textbf{Runner outcome:} \texttt{solved}\\
\textbf{ASP coverage:} positives \texttt{2/2}, negatives \texttt{1/1}

\subsubsection{Setup}

Same unit as cell 4: fork \(x_0\to x_1\), \(x_0\to x_2\) on the three-row source factorial, with asymmetric maps \(x_1:=2\) if \(x_0\neq 2\) else \(0\), and \(x_2:=2\) if \(x_0\neq 0\) else \(0\). For target \(x_2\), nonzero-positive labelling gives \(E^+=\{2,3\}\) and \(E^-=\{1\}\); on the table, \(x_2\neq 0\) if and only if \(x_0\neq 0\). Background knowledge includes exact-value predicates for the parent \(x_0\) and the sibling distractor \(x_1\); the target \(x_2\) is excluded. The sibling is correlated with \(x_2\) via \(x_0\) but imperfect as a separator (\texttt{x1\_val\_2} holds on both negative row 1 and positive row 2). The locked reference \(\mathcal{H}_{x_2}^\star\) is the two-clause parent expansion of \(x_0\neq 0\).

\subsubsection{Intended hypothesis \(\mathcal{H}_{x_2}^\star\)}

\begin{lstlisting}[style=logic,language=Prolog]
x2(A) :- x0_val_1(A).
x2(A) :- x0_val_2(A).
\end{lstlisting}

\subsubsection{What happened}

This cell uses the ECAI published configuration on unit U4 with learning target \(x_2\). After the initial optimum there are two rote rules, one for each positive. Non-deterministic folding replaces each identifier by the corresponding parent value in a single step, and both folds entail \(\langle E^+,E^-\rangle\) immediately. The learner returns a solved framework with full ASP coverage, no assumptions, and no citation of the sibling.

\subsubsection{What the learner did}

Learning begins from rote coverage of the two positives:

\begin{lstlisting}[style=logic,language=Prolog]
x2(A) :- A=2.
x2(A) :- A=3.
\end{lstlisting}

The first fold chosen is \(A=2\) (where \(x_0=1\)), which replaces the identifier by \texttt{x0\_val\_1}:

\begin{lstlisting}[style=logic,language=Prolog]
x2(A) :- x0_val_1(A).
\end{lstlisting}

That body does not cover the negative (\(x_0=0\)), so entailment succeeds without assumption introduction. The second fold is \(A=3\) (where \(x_0=2\)), which likewise yields a safe unary parent rule:

\begin{lstlisting}[style=logic,language=Prolog]
x2(A) :- x0_val_2(A).
\end{lstlisting}

Learning then stops with \texttt{nothing to fold}. Although sibling predicates (\texttt{x1\_val\_2}, \texttt{x1\_val\_0}) are present in the background knowledge, neither appears in a target body. Body-scope variables of the target rules are \(\{x_0\}\).

\subsubsection{Match assessment}

This is an exact match to \(\mathcal{H}_{x_2}^\star\).

The two unary parent rules are precisely the locked \texttt{val} expansion of ``\(x_0\neq 0\)'' under nonzero-positive labelling of \(x_2\). ASP coverage is complete, the solution uses no assumptions, and the sibling distractor is not cited.

\textbf{Comparison with cell 4.} On the same fork under ECAI, target \(x_1\) (cell 4) recovered \(\mathcal{H}_{x_1}^\star\) as \texttt{x0\_val\_0} / \texttt{x0\_val\_1} (parent \(x_0\neq 2\)), while this cell recovers \(\mathcal{H}_{x_2}^\star\) as \texttt{x0\_val\_1} / \texttt{x0\_val\_2} (parent \(x_0\neq 0\)). The learning path is the same in structure: two one-step parent folds, immediate entailment, no assumptions, sibling left unused. The difference is only which nonzero condition the asymmetric map induces, and which child appears in BK as the imperfect sibling distractor. Together, the two ECAI U4 cells show clean parent recovery on both branches of the fork.

\textbf{Short verdict.} Exact match to \(\mathcal{H}_{x_2}^\star\): two unary parent rules, no assumptions, sibling unused, full coverage (mirror of cell 4 on the other child).

\medskip\hrule\medskip

\subsection{6. ECAI \(\times\) \texttt{m12\_u5\_chain\_curated} \(\times\) \(t=x_2\)}

\textbf{Cell directory:} \texttt{M12x\_ecai2024/cells/m12\_u5\_chain\_curated\_\_target-x2/}\\
\textbf{Runner outcome:} \texttt{solved}\\
\textbf{ASP coverage:} positives \texttt{4/4}, negatives \texttt{2/2}

\subsubsection{Setup}

Unit U5 is the chain \(x_0\to x_1\to x_2\) on the curated six-row pilot support: for each \(x_0\in\{0,1,2\}\), the pairs \((x_0,x_0)\) and \((x_0,(x_0+1)\bmod 3)\), with \(x_2:=\operatorname{copy}(x_1)\). Nonzero-positive labelling gives \(E^+=\{2,3,4,5\}\) and \(E^-=\{1,6\}\); on the table, \(x_2\neq 0\) if and only if \(x_1\in\{1,2\}\). Background knowledge includes exact-value predicates for both ancestors \(x_0\) and \(x_1\); the target \(x_2\) is excluded. The direct parent \(x_1\) perfectly separates \(E^\pm\); the grandparent \(x_0\) is correlated but imperfect (\(x_0=0\) and \(x_0=2\) each contain one positive and one negative). The locked reference \(\mathcal{H}_{x_2}^\star\) prefers the parent copy expansion.

\subsubsection{Intended hypothesis \(\mathcal{H}_{x_2}^\star\)}

\begin{lstlisting}[style=logic,language=Prolog]
x2(A) :- x1_val_1(A).
x2(A) :- x1_val_2(A).
\end{lstlisting}

\subsubsection{What happened}

This cell uses the ECAI published configuration on unit U5 with learning target \(x_2\). After the initial optimum there are four rote rules, one for each positive. Non-deterministic folding proceeds by \(x_0\)-slices rather than by the parent \(x_1\): the complete \(x_0=1\) slice becomes a pure unary grandparent rule, while the incomplete \(x_0=0\) and \(x_0=2\) slices force assumption introduction. Contrary learning then recovers \(x_1\neq 0\) as the condition that defeats those assumptions. The learner returns a solved framework with full ASP coverage, but the target bodies cite the imperfect ancestor rather than the preferred parent.

\subsubsection{What the learner did}

Learning begins from rote coverage of the four positives. The first fold chosen is \(A=2\) (where \(x_0=0\), \(x_1=1\)), which replaces the identifier by \texttt{x0\_val\_0} alone:

\begin{lstlisting}[style=logic,language=Prolog]
x2(A) :- x0_val_0(A).
\end{lstlisting}

That body also covers the negative \(A=1\) (where \(x_0=0\), \(x_1=0\)), so entailment fails and assumption introduction adds \texttt{alpha\_1}:

\begin{lstlisting}[style=logic,language=Prolog]
x2(A) :- alpha_1(A), x0_val_0(A).
\end{lstlisting}

Next, \(A=3\) (where \(x_0=1\), \(x_1=1\)) folds to \texttt{x0\_val\_1}. Under this support every row with \(x_0=1\) is positive, so the unary grandparent rule entails immediately and needs no assumption; \(A=4\) is then deleted as subsumed:

\begin{lstlisting}[style=logic,language=Prolog]
x2(A) :- x0_val_1(A).
\end{lstlisting}

Finally \(A=5\) (where \(x_0=2\), \(x_1=2\)) folds to \texttt{x0\_val\_2}, which would also cover the negative \(A=6\), so assumption introduction adds \texttt{alpha\_2}:

\begin{lstlisting}[style=logic,language=Prolog]
x2(A) :- alpha_2(A), x0_val_2(A).
\end{lstlisting}

Contrary folding for \texttt{c\_alpha\_1} and \texttt{c\_alpha\_2} first tries the matching \(x_0\) literal (rejected: cannot introduce a further assumption on a body already tied to that assumption), then succeeds with the parent literal \texttt{x1\_val\_0} on both negatives:

\begin{lstlisting}[style=logic,language=Prolog]
c_alpha_1(A) :- x1_val_0(A).
c_alpha_2(A) :- x1_val_0(A).
\end{lstlisting}

In the ASP solution the assumptions are realised by

\begin{lstlisting}[style=logic,language=Prolog]
alpha_1(A) :- not c_alpha_1(A), x0_val_0(A).
alpha_2(A) :- not c_alpha_2(A), x0_val_2(A).
\end{lstlisting}

so on the incomplete \(x_0\) slices the assumption holds precisely when \(x_1\neq 0\). Body-scope variables of the target rules are \(\{x_0\}\); the preferred parent \(x_1\) appears only in the contrary conditions.

\subsubsection{Match assessment}

This is not an exact match to \(\mathcal{H}_{x_2}^\star\). The locked hypothesis is the unary parent expansion of \(x_1\in\{1,2\}\). The learner never emits \texttt{x1\_val\_1} or \texttt{x1\_val\_2} in a target body.

Extensionally the learned theory is an \(x_0\)-first case analysis repaired by parent-valued contraries:

\[ (x_0=0 \land x_1\neq 0)\;\lor\;(x_0=1)\;\lor\;(x_0=2 \land x_1\neq 0), \]

which coincides with \(x_1\neq 0\) on this curated table, so ASP coverage is complete. Intensionally this is the parent-vs-ancestor divergence the unit was designed to probe: the imperfect grandparent occupies the target bodies, and the true parent is recovered only as the contrary of the assumptions needed to repair overgeneral \(x_0\) folds.

The pattern is closely related to ECAI on U1 and U3: non-deterministic folding commits early to a one-literal body on a variable that does not perfectly separate \(E^\pm\) within its incomplete slices, then uses assumption/contrary machinery---here ultimately citing the preferred separator---to restore coverage. Relative to semantic success for U5, the framework covers the labelling but does not recover the intended parent-copy intension.

\textbf{Short verdict.} Not an exact match: grandparent \(x_0\) in target bodies with assumptions on incomplete slices; preferred parent \(x_1\) only in contraries; full coverage.

\medskip\hrule\medskip

\section{Batch B --- AAMAS (\texttt{M12x\_aamas2025})}

\subsection{10. AAMAS \(\times\) \texttt{m12\_u1\_separator\_copy} \(\times\) \(t=x_2\)}

\textbf{Cell directory:} \texttt{M12x\_aamas2025/cells/m12\_u1\_separator\_copy\_\_target-x2/}\\
\textbf{Runner outcome:} \texttt{solved}\\
\textbf{ASP coverage:} positives \texttt{6/6}, negatives \texttt{3/3}

\subsubsection{Setup}

Same unit and target as cell 1: separator \(x_1\to x_2\) with isolated \(x_0\);
\(x_2:=\operatorname{copy}(x_1)\) on the nine-row factorial; nonzero-positive examples
\(E^+=\{2,3,5,6,8,9\}\) and \(E^-=\{1,4,7\}\); exact-value background knowledge for
\(x_0\) and \(x_1\). This arm uses the AAMAS published configuration
(\texttt{folding\_mode(greedy)}, \texttt{folding\_selection(mgr)},
\texttt{folding\_space(bk)}, \texttt{asm\_intro(relto)}).

\subsubsection{Intended hypothesis \(\mathcal{H}_{x_2}^\star\)}

\begin{lstlisting}[style=logic,language=Prolog]
x2(A) :- x1_val_1(A).
x2(A) :- x1_val_2(A).
\end{lstlisting}

\subsubsection{What happened}

This cell uses the AAMAS published configuration on the same U1 instance. The learner
returned a solved framework with full ASP coverage and with no assumptions or
contraries. The learned target theory is a six-rule case split: each positive row
becomes a conjunction of that row's \(x_0\) value and \(x_1\) value. The parent \(x_1\)
appears in every body, but so does the isolated distractor \(x_0\). The result is
therefore not an exact match to \(\mathcal{H}_{x_2}^\star\).

\subsubsection{What the learner did}

Learning again begins from rote rules for the six positives. Under greedy folding with
selection \texttt{mgr} and folding space restricted to background knowledge, each
positive is folded by adjoining both column values present on that row. For example,
sample \(A=2\) has \(x_0=0\) and \(x_1=1\), and folds to:

\begin{lstlisting}[style=logic,language=Prolog]
x2(A) :- x0_val_0(A), x1_val_1(A).
\end{lstlisting}

Each such two-literal body entails the labelled example set immediately, so no
assumption introduction is triggered. The same pattern is applied independently to
every positive. The accepted target rules are:

\begin{lstlisting}[style=logic,language=Prolog]
x2(A) :- x0_val_0(A), x1_val_1(A).
x2(A) :- x0_val_0(A), x1_val_2(A).
x2(A) :- x0_val_1(A), x1_val_1(A).
x2(A) :- x0_val_1(A), x1_val_2(A).
x2(A) :- x0_val_2(A), x1_val_1(A).
x2(A) :- x0_val_2(A), x1_val_2(A).
\end{lstlisting}

There is no further generalisation that drops the \(x_0\) literals and merges the six
clauses into the two unary parent rules of \(\mathcal{H}_{x_2}^\star\). Body-scope
variables are \(\{x_0,x_1\}\). The solution uses no assumptions and no contraries.

\subsubsection{Match assessment}

The learned target rules are not an exact match to \(\mathcal{H}_{x_2}^\star\).

Extensionally, under mutual exclusivity and exhaustiveness of the three \(x_0\) value
predicates, the six conjunctions cover exactly the rows where \(x_1\in\{1,2\}\). That is
the same labelling as the intended two-rule parent hypothesis, which is why ASP coverage
is complete.

Intensionally the solution is a positive-row case split that retains the isolated
distractor in every body. The true parent is present, but the theory is a
parent-plus-distractor superset of \(\mathcal{H}_{x_2}^\star\), not the compact unary
parent rules. Relative to the ECAI twin on this unit, AAMAS avoids assumption-mediated
repair and keeps \(x_1\) in the target bodies, but it still fails to eliminate \(x_0\).

\textbf{Short verdict.} The cell is solved and extensionally adequate, with no assumptions,
but intensionally it enumerates six specialised conjunctions that keep the isolated
distractor beside the true parent.

\medskip\hrule\medskip

\subsection{11. AAMAS \(\times\) \texttt{m12\_u2\_collider\_min} \(\times\) \(t=x_2\)}

\textbf{Cell directory:} \texttt{M12x\_aamas2025/cells/m12\_u2\_collider\_min\_\_target-x2/}\\
\textbf{Runner outcome:} \texttt{solved}\\
\textbf{ASP coverage:} positives \texttt{4/4}, negatives \texttt{5/5}

\subsubsection{Setup}

Same unit and target as cell 2: collider \(x_0\to x_2\leftarrow x_1\) with \(x_2:=\min(x_0,x_1)\) on the nine-row factorial; nonzero-positive examples \(E^+=\{5,6,8,9\}\) and \(E^-=\{1,2,3,4,7\}\); exact-value background knowledge for both parents. This arm uses the AAMAS published configuration (\texttt{folding\_mode(greedy)}, \texttt{folding\_selection(mgr)}, \texttt{folding\_space(bk)}, \texttt{asm\_intro(relto)}).

\subsubsection{Intended hypothesis \(\mathcal{H}_{x_2}^\star\)}

\begin{lstlisting}[style=logic,language=Prolog]
x2(A) :- x0_val_1(A), x1_val_1(A).
x2(A) :- x0_val_1(A), x1_val_2(A).
x2(A) :- x0_val_2(A), x1_val_1(A).
x2(A) :- x0_val_2(A), x1_val_2(A).
\end{lstlisting}

\subsubsection{What happened}

This cell uses the AAMAS published configuration on the same U2 instance. After the initial optimum, the engine has four rote rules covering the positives by sample identifier:

\begin{lstlisting}[style=logic,language=Prolog]
x2(A) :- A=5.
x2(A) :- A=6.
x2(A) :- A=8.
x2(A) :- A=9.
\end{lstlisting}

The trace then computes a greedy folding plan for each of those four rules. For every positive, greedy/\texttt{mgr} first replaces the identifier by the \(x_0\) value on that row, then extends the body with the \(x_1\) value on the same row. The planned foldings are therefore the four nonzero parent pairs:

\begin{itemize}
  \item \(A=5\) \(\rightarrow\) \texttt{x0\_val\_1}, \texttt{x1\_val\_1}
  \item \(A=6\) \(\rightarrow\) \texttt{x0\_val\_1}, \texttt{x1\_val\_2}
  \item \(A=8\) \(\rightarrow\) \texttt{x0\_val\_2}, \texttt{x1\_val\_1}
  \item \(A=9\) \(\rightarrow\) \texttt{x0\_val\_2}, \texttt{x1\_val\_2}
\end{itemize}

Those foldings are then applied one at a time. After each application, entailment is checked against the \textbf{whole current theory}, not against the newly folded clause alone. So after the first fold, the theory still contains the remaining rote rules for \(A=6,8,9\); after the second, the remaining rotes for \(A=8,9\); and so on. Each intermediate theory already separates \(E^\pm\), so the trace reports \texttt{extended ABA entails <E+,E-> - using folding} and never introduces assumptions. Learning stops with \texttt{nothing to fold}. The final solution is the four specialised conjunctions, with no assumptions and no contraries, and with full ASP coverage.

\subsubsection{What the learner did}

The accepted target rules are:

\begin{lstlisting}[style=logic,language=Prolog]
x2(A) :- x0_val_1(A), x1_val_1(A).   % from A=5
x2(A) :- x0_val_1(A), x1_val_2(A).   % from A=6
x2(A) :- x0_val_2(A), x1_val_1(A).   % from A=8
x2(A) :- x0_val_2(A), x1_val_2(A).   % from A=9
\end{lstlisting}

Body-scope variables are \(\{x_0,x_1\}\).

\subsubsection{Match assessment}

This is an exact match to \(\mathcal{H}_{x_2}^\star\).

The four conjunctions are precisely the locked \texttt{val} expansion of ``both parents are nonzero.'' ASP coverage is complete, and the solution uses no assumptions.

\textbf{Contrast with ECAI on the same unit.} Under ECAI (cell 2), non-deterministic folding first produces one-literal bodies such as \texttt{x2(A) :- x0\_val\_1(A)}, which overgeneralise within an \(x_0\)-slice and force assumption introduction; the second parent is recovered only through contraries, yielding a compact assumption-mediated encoding of both-parents-nonzero that is semantically successful but not an exact syntactic match to \(\mathcal{H}_{x_2}^\star\). Under AAMAS, greedy folding specialises each positive to the full parent pair immediately, so each intermediate theory remains safe without assumptions and the final rules coincide with \(\mathcal{H}_{x_2}^\star\). Both arms recover the intended concept; they differ in presentation: ECAI's assumption form versus AAMAS's exact case enumeration.

\textbf{Short verdict.} Exact match to \(\mathcal{H}_{x_2}^\star\): four both-parent conjunctions, no assumptions, full coverage.

\medskip\hrule\medskip

\subsection{12. AAMAS \(\times\) \texttt{m12\_u3\_collider\_max} \(\times\) \(t=x_2\)}

\textbf{Cell directory:} \texttt{M12x\_aamas2025/cells/m12\_u3\_collider\_max\_\_target-x2/}\\
\textbf{Runner outcome:} \texttt{solved}\\
\textbf{ASP coverage:} positives \texttt{8/8}, negatives \texttt{1/1}

\subsubsection{Setup}

Same unit and target as cell 3: collider \(x_0\to x_2\leftarrow x_1\) with \(x_2:=\max(x_0,x_1)\) on the nine-row factorial; under nonzero-positive labelling, \(x_2\neq 0\) if and only if at least one parent is nonzero. The positive examples are rows \(\{2,3,4,5,6,7,8,9\}\) and the sole negative is row \(\{1\}\). Background knowledge contains exact-value predicates for both parents; the target \(x_2\) is excluded. This arm uses the AAMAS published configuration (\texttt{folding\_mode(greedy)}, \texttt{folding\_selection(mgr)}, \texttt{folding\_space(bk)}, \texttt{asm\_intro(relto)}). The locked reference \(\mathcal{H}_{x_2}^\star\) is the four-clause \texttt{val} expansion of ``at least one parent is nonzero.''

\subsubsection{Intended hypothesis \(\mathcal{H}_{x_2}^\star\)}

\begin{lstlisting}[style=logic,language=Prolog]
x2(A) :- x0_val_1(A).
x2(A) :- x0_val_2(A).
x2(A) :- x1_val_1(A).
x2(A) :- x1_val_2(A).
\end{lstlisting}

\subsubsection{What happened}

This cell uses the AAMAS published configuration on the same U3 instance. After the initial optimum, the engine has eight rote rules covering the positives by sample identifier. The trace then computes a greedy folding plan for each of those rules. For every positive, greedy/\texttt{mgr} first replaces the identifier by the \(x_0\) value on that row, then extends the body with the \(x_1\) value on the same row. The planned foldings are therefore the eight nonzero parent pairs of the factorial (every cell except \((x_0,x_1)=(0,0)\)):

\begin{itemize}
  \item \(A=2\) \(\rightarrow\) \texttt{x0\_val\_0}, \texttt{x1\_val\_1}
  \item \(A=3\) \(\rightarrow\) \texttt{x0\_val\_0}, \texttt{x1\_val\_2}
  \item \(A=4\) \(\rightarrow\) \texttt{x0\_val\_1}, \texttt{x1\_val\_0}
  \item \(A=5\) \(\rightarrow\) \texttt{x0\_val\_1}, \texttt{x1\_val\_1}
  \item \(A=6\) \(\rightarrow\) \texttt{x0\_val\_1}, \texttt{x1\_val\_2}
  \item \(A=7\) \(\rightarrow\) \texttt{x0\_val\_2}, \texttt{x1\_val\_0}
  \item \(A=8\) \(\rightarrow\) \texttt{x0\_val\_2}, \texttt{x1\_val\_1}
  \item \(A=9\) \(\rightarrow\) \texttt{x0\_val\_2}, \texttt{x1\_val\_2}
\end{itemize}

Those foldings are then applied one at a time. After each application, entailment is checked against the \textbf{whole current theory}, not against the newly folded clause alone. Each pairwise body matches only its own factorial cell, so it never covers the negative \(A=1\). Every intermediate theory already separates \(E^\pm\), the trace reports \texttt{extended ABA entails <E+,E-> - using folding}, and no assumptions are introduced. Learning stops with \texttt{nothing to fold}. The final solution is the eight specialised conjunctions, with no assumptions and no contraries, and with full ASP coverage.

\subsubsection{What the learner did}

The accepted target rules are:

\begin{lstlisting}[style=logic,language=Prolog]
x2(A) :- x0_val_0(A), x1_val_1(A).
x2(A) :- x0_val_0(A), x1_val_2(A).
x2(A) :- x0_val_1(A), x1_val_0(A).
x2(A) :- x0_val_1(A), x1_val_1(A).
x2(A) :- x0_val_1(A), x1_val_2(A).
x2(A) :- x0_val_2(A), x1_val_0(A).
x2(A) :- x0_val_2(A), x1_val_1(A).
x2(A) :- x0_val_2(A), x1_val_2(A).
\end{lstlisting}

Body-scope variables are \(\{x_0,x_1\}\).

\subsubsection{Match assessment}

This is not an exact match to \(\mathcal{H}_{x_2}^\star\). The locked hypothesis is the compact four-rule OR of unary nonzero conditions on either parent. The learner instead emits an exhaustive enumeration of the eight positive cells.

Extensionally the two presentations coincide: the eight conjunctions cover precisely those rows where \(\max(x_0,x_1)\neq 0\), i.e. at least one parent is nonzero. ASP coverage is complete for that reason. Intensionally, however, the solution never discovers a one-literal generalisation such as \texttt{x2(A) :- x0\_val\_1(A)}; under greedy/\texttt{mgr} each positive is specialised to the full parent pair immediately, and that pair already entails safely, so there is no pressure to fold further or to introduce assumptions.

\textbf{Contrast with ECAI on the same unit.} Under ECAI (cell 3), non-deterministic folding produces an \(x_0\)-first case analysis: pure unary rules on the complete \(x_0\in\{1,2\}\) slices, and an assumption-mediated rule only on the incomplete \(x_0=0\) slice. That is more compact than eight pairwise clauses and closer in spirit to a disjunctive reading, though still not the locked four-rule OR. Under AAMAS, the same unit yields the exhaustive positive-cell enumeration with both parents in every body and no assumptions. Both arms recover the intended concept; they differ in presentation: ECAI's mixed unary/assumption case analysis versus AAMAS's eight specialised conjunctions.

\textbf{Contrast with AAMAS on U2.} The same greedy pairwise specialisation matched \(\mathcal{H}^\star\) exactly on U2, because the locked min hypothesis \emph{is} a set of both-parent conjunctions. On U3 the locked max hypothesis is a unary OR, so the same folding style recovers the concept extensionally but misses \(\mathcal{H}_{x_2}^\star\) syntactically.

\textbf{Short verdict.} Not an exact match to the locked four-rule OR: eight both-parent conjunctions enumerating the positive cells, no assumptions, full coverage of the intended ``at least one parent nonzero'' concept.

\medskip\hrule\medskip

\subsection{13. AAMAS \(\times\) \texttt{m12\_u4\_fork\_asymmetric} \(\times\) \(t=x_1\)}

\textbf{Cell directory:} \texttt{M12x\_aamas2025/cells/m12\_u4\_fork\_asymmetric\_\_target-x1/}\\
\textbf{Runner outcome:} \texttt{solved}\\
\textbf{ASP coverage:} positives \texttt{2/2}, negatives \texttt{1/1}

\subsubsection{Setup}

Same unit and target as cell 4: fork \(x_0\to x_1\), \(x_0\to x_2\) on the three-row source factorial; \(x_1:=2\) if \(x_0\neq 2\) else \(0\); nonzero-positive examples \(E^+=\{1,2\}\) and \(E^-=\{3\}\); on the table, \(x_1\neq 0\) if and only if \(x_0\in\{0,1\}\). Background knowledge includes exact-value predicates for the parent \(x_0\) and the sibling distractor \(x_2\). This arm uses the AAMAS published configuration (\texttt{folding\_mode(greedy)}, \texttt{folding\_selection(mgr)}, \texttt{folding\_space(bk)}, \texttt{asm\_intro(relto)}). The locked reference \(\mathcal{H}_{x_1}^\star\) is the two-clause unary parent expansion of \(x_0\neq 2\).

\subsubsection{Intended hypothesis \(\mathcal{H}_{x_1}^\star\)}

\begin{lstlisting}[style=logic,language=Prolog]
x1(A) :- x0_val_0(A).
x1(A) :- x0_val_1(A).
\end{lstlisting}

\subsubsection{What happened}

This cell uses the AAMAS published configuration on the same U4 \(x_1\) instance. After the initial optimum, the engine has two rote rules covering the positives by sample identifier. The trace then computes a greedy folding plan for each rule. For every positive, greedy/\texttt{mgr} first replaces the identifier by the \(x_0\) value on that row, then extends the body with the sibling \(x_2\) value on the same row. The planned foldings are therefore the observed parent--sibling pairs:

\begin{itemize}
  \item \(A=1\) \(\rightarrow\) \texttt{x0\_val\_0}, \texttt{x2\_val\_0}
  \item \(A=2\) \(\rightarrow\) \texttt{x0\_val\_1}, \texttt{x2\_val\_2}
\end{itemize}

Those foldings are applied one at a time. Each pairwise body matches only its own row on this three-row table, so neither covers the negative \(A=3\). Every intermediate theory already separates \(E^\pm\), the trace reports \texttt{extended ABA entails <E+,E-> - using folding}, and no assumptions are introduced. Learning stops with \texttt{nothing to fold}. The final solution cites both the true parent and the sibling distractor in every target body.

\subsubsection{What the learner did}

The accepted target rules are:

\begin{lstlisting}[style=logic,language=Prolog]
x1(A) :- x0_val_0(A), x2_val_0(A).
x1(A) :- x0_val_1(A), x2_val_2(A).
\end{lstlisting}

Body-scope variables are \(\{x_0,x_2\}\).

\subsubsection{Match assessment}

This is not an exact match to \(\mathcal{H}_{x_1}^\star\). The locked hypothesis uses unary parent conditions only. The learner instead emits two specialised conjunctions that retain the sibling value observed on each positive row.

Extensionally the theory still separates \(E^\pm\) with full ASP coverage: on this table the parent literals alone already suffice, and the sibling conjuncts are redundant specialisations. Intensionally this is a sibling-citation divergence. Relative to the unit's semantic success description, the parent is recovered, but the imperfect sibling distractor is also cited---the failure mode the U4 design was meant to probe.

\textbf{Comparison with ECAI on the same cell (cell 4).} Under ECAI, non-deterministic folding stops after the one-literal parent folds \texttt{x0\_val\_0} and \texttt{x0\_val\_1}, which already entail safely, yielding an exact match to \(\mathcal{H}_{x_1}^\star\) with the sibling unused. Under AAMAS, greedy/\texttt{mgr} continues from the parent literal to the sibling literal on the same row before accepting the fold, so the sibling enters every body even though it is not needed for entailment. Both arms recover a covering theory centred on the parent; they differ in whether the correlated sibling is left out of the intension.

\textbf{Short verdict.} Not an exact match: parent recovered, but sibling \(x_2\) cited in both bodies via greedy pairwise specialisation; full coverage, no assumptions.

\medskip\hrule\medskip

\subsection{14. AAMAS \(\times\) \texttt{m12\_u4\_fork\_asymmetric} \(\times\) \(t=x_2\)}

\textbf{Cell directory:} \texttt{M12x\_aamas2025/cells/m12\_u4\_fork\_asymmetric\_\_target-x2/}\\
\textbf{Runner outcome:} \texttt{solved}\\
\textbf{ASP coverage:} positives \texttt{2/2}, negatives \texttt{1/1}

\subsubsection{Setup}

Same unit and target as cell 5: fork \(x_0\to x_1\), \(x_0\to x_2\) on the three-row source factorial; \(x_2:=2\) if \(x_0\neq 0\) else \(0\); nonzero-positive examples \(E^+=\{2,3\}\) and \(E^-=\{1\}\); on the table, \(x_2\neq 0\) if and only if \(x_0\neq 0\). Background knowledge includes exact-value predicates for the parent \(x_0\) and the sibling distractor \(x_1\). This arm uses the AAMAS published configuration (\texttt{folding\_mode(greedy)}, \texttt{folding\_selection(mgr)}, \texttt{folding\_space(bk)}, \texttt{asm\_intro(relto)}). The locked reference \(\mathcal{H}_{x_2}^\star\) is the two-clause unary parent expansion of \(x_0\neq 0\).

\subsubsection{Intended hypothesis \(\mathcal{H}_{x_2}^\star\)}

\begin{lstlisting}[style=logic,language=Prolog]
x2(A) :- x0_val_1(A).
x2(A) :- x0_val_2(A).
\end{lstlisting}

\subsubsection{What happened}

This cell uses the AAMAS published configuration on the same U4 \(x_2\) instance. After the initial optimum, the engine has two rote rules covering the positives by sample identifier. The trace then computes a greedy folding plan for each rule. For every positive, greedy/\texttt{mgr} first replaces the identifier by the \(x_0\) value on that row, then extends the body with the sibling \(x_1\) value on the same row. The planned foldings are therefore the observed parent--sibling pairs:

\begin{itemize}
  \item \(A=2\) \(\rightarrow\) \texttt{x0\_val\_1}, \texttt{x1\_val\_2}
  \item \(A=3\) \(\rightarrow\) \texttt{x0\_val\_2}, \texttt{x1\_val\_0}
\end{itemize}

Those foldings are applied one at a time. Each pairwise body matches only its own row on this three-row table, so neither covers the negative \(A=1\). Every intermediate theory already separates \(E^\pm\), the trace reports \texttt{extended ABA entails <E+,E-> - using folding}, and no assumptions are introduced. Learning stops with \texttt{nothing to fold}. The final solution cites both the true parent and the sibling distractor in every target body.

\subsubsection{What the learner did}

The accepted target rules are:

\begin{lstlisting}[style=logic,language=Prolog]
x2(A) :- x0_val_1(A), x1_val_2(A).
x2(A) :- x0_val_2(A), x1_val_0(A).
\end{lstlisting}

Body-scope variables are \(\{x_0,x_1\}\).

\subsubsection{Match assessment}

This is not an exact match to \(\mathcal{H}_{x_2}^\star\). The locked hypothesis uses unary parent conditions only. The learner instead emits two specialised conjunctions that retain the sibling value observed on each positive row.

Extensionally the theory still separates \(E^\pm\) with full ASP coverage: on this table the parent literals alone already suffice, and the sibling conjuncts are redundant specialisations. Intensionally this is a sibling-citation divergence of the same kind as on AAMAS \(x_1\).

\textbf{Comparison with ECAI on the same cell (cell 5).} Under ECAI, non-deterministic folding stops after the one-literal parent folds \texttt{x0\_val\_1} and \texttt{x0\_val\_2}, yielding an exact match to \(\mathcal{H}_{x_2}^\star\) with the sibling unused. Under AAMAS, greedy/\texttt{mgr} continues to the sibling literal on the same row, so \(x_1\) enters every body even though it is not needed for entailment.

\textbf{Comparison with AAMAS on \(x_1\) (cell 13).} The learning path is the same on both branches of the fork: parent first, then sibling, accept pairwise specialisations without assumptions. Cell 13 cites distractor \(x_2\) beside parent \(x_0\); this cell cites distractor \(x_1\) beside parent \(x_0\). Together with the two exact ECAI U4 cells, the four U4 outcomes show a clean arm contrast: ECAI recovers unary parent \(\mathcal{H}^\star\) on both children; AAMAS recovers covering parent-centred rules that also cite the correlated sibling.

\textbf{Short verdict.} Not an exact match: parent recovered, but sibling \(x_1\) cited in both bodies via greedy pairwise specialisation; full coverage, no assumptions (mirror of cell 13).

\medskip\hrule\medskip

\subsection{15. AAMAS \(\times\) \texttt{m12\_u5\_chain\_curated} \(\times\) \(t=x_2\)}

\textbf{Cell directory:} \texttt{M12x\_aamas2025/cells/m12\_u5\_chain\_curated\_\_target-x2/}\\
\textbf{Runner outcome:} \texttt{solved}\\
\textbf{ASP coverage:} positives \texttt{4/4}, negatives \texttt{2/2}

\subsubsection{Setup}

Same unit and target as cell 6: chain \(x_0\to x_1\to x_2\) on the curated six-row support with \(x_2:=\operatorname{copy}(x_1)\); nonzero-positive examples \(E^+=\{2,3,4,5\}\) and \(E^-=\{1,6\}\); on the table, \(x_2\neq 0\) if and only if \(x_1\in\{1,2\}\). Background knowledge includes exact-value predicates for both ancestors. This arm uses the AAMAS published configuration (\texttt{folding\_mode(greedy)}, \texttt{folding\_selection(mgr)}, \texttt{folding\_space(bk)}, \texttt{asm\_intro(relto)}). The locked reference \(\mathcal{H}_{x_2}^\star\) prefers the unary parent copy expansion.

\subsubsection{Intended hypothesis \(\mathcal{H}_{x_2}^\star\)}

\begin{lstlisting}[style=logic,language=Prolog]
x2(A) :- x1_val_1(A).
x2(A) :- x1_val_2(A).
\end{lstlisting}

\subsubsection{What happened}

This cell uses the AAMAS published configuration on the same U5 instance. After the initial optimum, the engine has four rote rules covering the positives by sample identifier. The trace then computes a greedy folding plan for each rule. For every positive, greedy/\texttt{mgr} first replaces the identifier by the grandparent \(x_0\) value on that row, then extends the body with the parent \(x_1\) value on the same row. The planned foldings are therefore the four observed ancestor--parent pairs on the positives:

\begin{itemize}
  \item \(A=2\) \(\rightarrow\) \texttt{x0\_val\_0}, \texttt{x1\_val\_1}
  \item \(A=3\) \(\rightarrow\) \texttt{x0\_val\_1}, \texttt{x1\_val\_1}
  \item \(A=4\) \(\rightarrow\) \texttt{x0\_val\_1}, \texttt{x1\_val\_2}
  \item \(A=5\) \(\rightarrow\) \texttt{x0\_val\_2}, \texttt{x1\_val\_2}
\end{itemize}

Those foldings are applied one at a time. Each pairwise body matches only its own row on this support, so none covers a negative. Every intermediate theory already separates \(E^\pm\), the trace reports \texttt{extended ABA entails <E+,E-> - using folding}, and no assumptions are introduced. Learning stops with \texttt{nothing to fold}. The final solution cites both the preferred parent and the imperfect grandparent in every target body.

\subsubsection{What the learner did}

The accepted target rules are:

\begin{lstlisting}[style=logic,language=Prolog]
x2(A) :- x0_val_0(A), x1_val_1(A).
x2(A) :- x0_val_1(A), x1_val_1(A).
x2(A) :- x0_val_1(A), x1_val_2(A).
x2(A) :- x0_val_2(A), x1_val_2(A).
\end{lstlisting}

Body-scope variables are \(\{x_0,x_1\}\).

\subsubsection{Match assessment}

This is not an exact match to \(\mathcal{H}_{x_2}^\star\). The locked hypothesis is the compact unary parent expansion. The learner instead emits four specialised conjunctions that retain the grandparent value observed on each positive row.

Extensionally the theory separates \(E^\pm\) with full ASP coverage: each rule is a positive cell of the curated support, and together they cover exactly the rows where \(x_1\in\{1,2\}\). Intensionally the parent \(x_1\) is present in every body---so the preferred separator is recovered as a conjunct---but the imperfect ancestor \(x_0\) is also cited throughout. That is the parent-vs-ancestor contamination the unit watches for: success is not the clean parent-only intension of \(\mathcal{H}_{x_2}^\star\).

\textbf{Comparison with ECAI on the same cell (cell 6).} Under ECAI, non-deterministic folding commits to unary \(x_0\) bodies, forces assumptions on the incomplete \(x_0\) slices, and recovers \(x_1\neq 0\) only through contraries; the preferred parent never appears in a target body. Under AAMAS, greedy specialisation keeps both ancestors in every body, needs no assumptions, and does place the parent in the target rules---but only beside the grandparent. Both arms cover the labelling and both miss exact \(\mathcal{H}_{x_2}^\star\); they differ in where the parent appears (contrary-only under ECAI versus redundant conjunct under AAMAS) and in whether assumptions are used.

\textbf{Short verdict.} Not an exact match: four parent--grandparent conjunctions; parent present but ancestor retained; no assumptions, full coverage.

\medskip\hrule\medskip

\emph{Cells 7--9 and 16--18 to follow.}

\end{document}
